\addcontentsline{toc}{chapter}{Abstract}
\begin{center}
{\bf A B S T R A C T}
\end{center}

\setlength{\parindent}{1.5em}
\sloppy

Sanskrit, with its formal grammatical foundations and rich morphological structure, presents both an opportunity and a challenge for computational modeling. Rooted in the generative principles of Pāṇini's \textit{Aṣṭādhyāyī}, the language demands a parsing approach that can navigate complex linguistic phenomena such as \textit{sandhi}, compounding, and free word order. Traditional rule-based systems offer structural rigor but struggle with adaptability, while neural models excel in generalization but often lack linguistic grounding.

This thesis proposes a hybrid framework that seeks to integrate grammatical constraints into neural architectures for Sanskrit parsing. The envisioned system emphasizes sandhi awareness and aims to balance symbolic precision with the flexibility of data-driven models. While implementation and evaluation are planned in future stages, the current work establishes the theoretical foundation and design direction. Broadly, this research aspires to contribute toward linguistically informed NLP for Sanskrit and similar morphologically rich, low-resource languages by exploring the intersection of classical grammar and modern computation.
